% Complete proof candidate. Individual check status is in result_audit.json.
\subsection{C.6. Supplement: exact leading factors from a genus-character denominator bound}

We retain the full real period $\Omega_i$ and the absolute canonical
height $\widehat h_O$ of the preceding subsection. The argument below
uses the Gross--Zagier formula to bound a rational denominator; it does
not assume the leading-term Birch--Swinnerton-Dyer formula.

\begin{proposition}[ARI C.6.R5; Proved: all remaining rational leading factors]\label{ari:C.6.R5}
For the six rank-one rows $i=2,4,6,9,10,11$, one has
\[
 \frac{L'(E_i,1)}{\Omega_i}
   =q_i\widehat h_O(P_i),\qquad
 (q_2,q_4,q_6,q_9,q_{10},q_{11})=(8,32,8,16,16,576).
\]
Together with C.6.R4, this determines every rank-one factor:
\[
\begin{array}{c|rrrrrrrrr}
 i&2&3&4&5&6&7&9&10&11\\\hline
 q_i&8&4&32&4&8&4&16&16&576.
\end{array}
\]
The certified auxiliary ratios and denominator bounds used to identify
the six additional factors are as follows. Here $D$ is the auxiliary
negative fundamental discriminant,
$r_i=L(E_i^D,1)/\Omega_{i,D}$, and the denominator of $q_i$ is at most
$B_i$.
\[
\begin{array}{c|r|r|r|r|r}
 i&D&r_i&B_i&q_i&|a_i|\\\hline
 2&-23&64&256&8&32\\
 4&-31&64&256&32&64\\
 6&-19&8&32&8&8\\
 9&-19&16&64&16&16\\
 10&-3&400&1600&16&80\\
 11&-3&64&256&576&192
\end{array}
\]
The last column is the absolute genus-trace integer in the proof. These
identities determine the analytic BSD quotients exactly: they are one
for $i=2,4,6,9,10$ and nine for $i=11$. They do not by themselves
identify those numbers with the orders of the Tate--Shafarevich groups.
\end{proposition}

\begin{proof}
We give the arithmetic denominator bound first and then the interval
certificate that isolates the listed rational numbers.

\paragraph{1. A parametrization under quadratic twisting.}
Let $E_0/\Q$ have conductor $N_0$, and suppose a parametrization
$\phi:X_0(N_0)\to E_0$ sends $\infty$ to $O$ and satisfies
\[
 \phi^*\omega_0=c_0f_0(q)\,dq/q,
\]
where $\omega_0$ is a minimal differential and $f_0$ is the normalized
newform. Let $b$ be an odd fundamental discriminant coprime to $N_0$,
and suppose that the minimal twist has conductor $N_0b^2$ and exact
invariants
\[
 c_4(E_0^b)=b^2c_4(E_0),\qquad
 c_6(E_0^b)=b^3c_6(E_0).
\]
Then $E_0^b$ has a parametrization from $X_0(N_0b^2)$ with the same
Manin factor $c_0$.

Indeed, put $m=|b|$ and, using addition on $E_0$, define
\[
 F(z)=\sum_{a\bmod m}\chi_b(a)\phi(z+a/m).
\]
For $\gamma=\left(\begin{smallmatrix}A&B\\C&D_0\end{smallmatrix}\right)$
in $\Gamma_0(N_0m^2)$ choose $a'\equiv aD_0^2\pmod m$.
Direct multiplication gives
$T_{a/m}\gamma T_{a'/m}^{-1}\in\Gamma_0(N_0)$, where $T_t$
denotes translation by $t$. Because $\chi_b(D_0^2)=1$, the sum is
invariant. Each summand is an algebraic morphism on a common
congruence cover; invariance therefore descends the sum to
$X_0(N_0m^2)$ over $\Q(\zeta_m)$. Its expansion at $\infty$
gives $F^{\sigma_u}=\chi_b(u)F$ for
$\sigma_u(\zeta_m)=\zeta_m^u$. Quadratic descent gives the required
map to $E_0^b$ over $\Q$, with $\infty$ mapping to $O$.
On differentials, the Gauss-sum identity gives
\[
 F^*\omega_0=c_0\tau(\chi_b)f_b(q)\,dq/q,
 \qquad \tau(\chi_b)^2=b.
\]
For the twist isomorphism $\iota:E_0^b\to E_0$ the exact invariant
identities imply
$\iota^*\omega_0=\sqrt b\,\omega_b$ up to sign, with no additional
minimal-model scaling. Cancellation proves the assertion. The map
need not be optimal; optimality is not used below.

If $c_0=1$, let $t_b=\#E_0^b(\Q)_{\rm tors}$ and
$n_b=\#\pi_0(E_0^b(\R))$. Manin--Drinfeld says that the image of
the rational cuspidal divisor $[0]-[\infty]$ is torsion. Its elliptic
logarithm, along the imaginary axis, is $\pm L(E_0^b,1)$.
Consequently $t_bL(E_0^b,1)$ is an integral multiple of the positive
real lattice period $\Omega_b^+$, and
\begin{equation}\label{ari:eq:Cgz-central-integrality}
 t_bn_b\,\frac{L(E_0^b,1)}{\Omega_b}\in\Z,
 \qquad \Omega_b=n_b\Omega_b^+.
\end{equation}

\paragraph{2. The six small base curves and Heegner hypotheses.}
Use the following data; $E_i=E_0^d$ and $K=\Q(\sqrt D)$.
\[
\begin{array}{c|c|r|r|r}
 i&E_0&N_0&d&D\\\hline
 2&576\mathrm e2&576&-11&-23\\
 4&3136\mathrm r2&3136&57&-31\\
 6&121\mathrm b1&121&21&-19\\
 9&1849\mathrm a1&1849&21&-19\\
 10&4489\mathrm a1&4489&217&-3\\
 11&26569\mathrm a1&26569&185801&-3
\end{array}
\]
All six base Manin factors are one. For the first five bases this
is independently returned by PARI's unconditional modular-symbol
algorithm. For $26569\mathrm a1$, the value one and modular degree
$59368$ are explicitly tabulated by
\citet[Appendix B, p.~102, unnumbered $26569\mathrm a1/26569\mathrm a2$
block]{CgzBelabas2019}; their Algorithm 1 describes the computation.
The first five completed computations are in
\path{code/B_C6_gz_setup.out}. The separate exact audit
\path{code/B_C6_gz_exact_metadata.gp} checks the minimal equations,
both $c_4,c_6$ scaling identities for $b=d,dD$, both conductor
identities, all torsion orders, and all splitting conditions. It
does not recognize any of these identities from floating-point
scaling factors.

In every row $d,D,dD$ are fundamental discriminants,
$\gcd(d,D)=\gcd(dD,N_0)=1$, and every prime dividing $N_0$ splits
in $K$. The genus character for $K(\sqrt d)/K$ factors through
$\operatorname{Pic}(\mathcal O_{|d|})$ and has conductor exactly
$|d|$. On ideals prime to $d$ it is
\[
 \chi(\mathfrak a)=\chi_d(N\mathfrak a).
\]
For a unit of the order locally at a prime $p\mid d$, its norm is
a square modulo $p$, so this character is trivial on the order's
units. It is nontrivial on the full local unit group: at a split
prime the norm is multiplication on two copies of $\Q_p^\times$,
and at an inert prime the residue-field norm onto $\mathbb F_p^\times$
is surjective. Since $p$ is odd and unramified in $K$, this also
proves that the conductor contains every $p\mid d$ exactly once.
Thus the primitive ring-class-character hypotheses of
\citet[Theorem 1.1, pp.~2524--2525]{CgzCai2014} hold. Their
$\mu(N_0,D)=0$ and
$u=[\mathcal O_{|d|}^{\times}:\Z^{\times}]=1$; the latter remains
true for $D=-3$ because $|d|>1$.

Let $P_{\mathfrak a}$ be the corresponding Heegner points and put
\[
 Y=\sum_{[\mathfrak a]}\chi(\mathfrak a)\phi(P_{\mathfrak a}).
\]
The coefficients are $\pm1$, so this is an actual algebraic point,
not merely an element of a tensor product with $\Q$. It is fixed
by $\ker\chi$ and anti-invariant under the other automorphism of
$K(\sqrt d)/K$. Twisting therefore gives $Y_d\in E_i(K)$.
Artin factorization gives
\[
 L(s,E_0,\chi)=L(E_i,s)L(E_i^D,s).
\]
The intervals below prove $L(E_i^D,1)>0$, hence the auxiliary
curve has rank zero. The rank-zero implication used here is the
same unconditional modular/CM rank implication used in C.6.R2--R3.
Thus $E_i(K)$ has rank one. C.6.R1 proves that $P_i$ is primitive
over $\Q$. Since $Y_d-\overline{Y_d}$ is torsion and
$Y_d+\overline{Y_d}$ is rational, there is an integer $a_i$ with
\[
 2Y_d=a_iP_i\pmod{E_i(K)_{\rm tors}},\qquad
 \widehat h_O(Y_d)=\frac{a_i^2}{4}\widehat h_O(P_i).
\]

\paragraph{3. The exact height and period normalization.}
Write $V_0=\operatorname{area}(\C/\Lambda_0)$ for the period area
of the minimal differential. Integrating $\phi^*\omega_0$ gives
\[
 4\pi^2(f_0,f_0)_{\Gamma_0(N_0)}
   =\frac{\deg\phi}{c_0^2}V_0.
\]
The degree in the denominator of CST's Theorem 1.1 cancels, giving
\begin{equation}\label{ari:eq:Cgz-product}
 L'(E_i,1)L(E_i^D,1)
  =\frac{2V_0}{c_0^2|d|\sqrt{|D|}}\,\widehat h_K(Y).
\end{equation}
Here $\widehat h_K$ is the diagonal N\'eron--Tate pairing relative
to $K$. This convention is explicit in
\citet[\S7.1.1, pp.~218--219, Theorem 7.1.2]{CgzYZZ2011}:
their underlying height is normalized by $[L:K]^{-1}$ over a
field of definition $L$, and their pairing is
$\widehat h_{\mathcal L}(x+y)-\widehat h_{\mathcal L}(x)
-\widehat h_{\mathcal L}(y)$.

One can also see the factor geometrically. For
$\mathcal L=\mathcal O(O)$, pull the normalized Poincar\'e bundle
back under $1\times\lambda_{\mathcal L}$. The result is
$m^*\mathcal L\otimes p_1^*\mathcal L^{-1}\otimes
p_2^*\mathcal L^{-1}$. On the diagonal its canonical height is
that of $[2]^*\mathcal L\otimes\mathcal L^{-2}\simeq
\mathcal L^2$, namely $2\widehat h_O$ in the absolute convention.
Changing from the absolute convention to relative heights over
$K$ multiplies by $[K:\Q]=2$. Therefore
\[
 \widehat h_K(Y)=4\widehat h_O(Y_d)
                =a_i^2\widehat h_O(P_i).
\]
This factor comes from the Poincar\'e pairing and the base field,
independently of any numerical height routine.

A real-adapted basis of $\Lambda_0$ gives
\[
 \Omega_i\Omega_{i,D}=
 \begin{cases}
 2V_0/(|d|\sqrt{|D|}),&\Delta(E_0)<0,\\
 4V_0/(|d|\sqrt{|D|}),&\Delta(E_0)>0.
 \end{cases}
\]
Indeed, a triangular real lattice has one real component and its
purely imaginary primitive period has twice the height of its
nonreal basis vector. A rectangular lattice has two real
components. Twisting by the two opposite signs interchanges the
real and imaginary directions, and the exact minimal differential
scalings in step 2 supply $1/(|d|\sqrt{|D|})$.
Substituting into \eqref{ari:eq:Cgz-product} gives
\begin{equation}\label{ari:eq:Cgz-denominator}
 q_ir_i=
 \begin{cases}
 a_i^2/c_0^2,&\Delta(E_0)<0,\\
 a_i^2/(2c_0^2),&\Delta(E_0)>0.
 \end{cases}
\end{equation}
In particular, if $r_i=A_i/B_i'$ is reduced and $c_0=1$, the
denominator of $q_i$ divides $2|A_i|$. We use the weaker upper
bound $B_i=4|A_i|$ in the interval certificate.

\paragraph{4. Global identification of the CM coefficient series.}
For $s=11,43,67,163$, put
$F_s=\Q(\sqrt{-s})$ and $\mathfrak p=(\sqrt{-s})$.
These fields have class number one and units $\{\pm1\}$.
The character
\[
 \psi((\alpha))=
  \left(\frac{\alpha\bmod\mathfrak p}{s}\right)\alpha
\]
is well defined because $(-1/s)=-1$. It has infinity type $(1,0)$
and finite conductor $\mathfrak p$. Hecke's theta construction
places $\sum_{\mathfrak a}\psi(\mathfrak a)q^{N\mathfrak a}$ in
$S_2(\Gamma_0(s^2))$: the nebentypus is
$\chi_{-s}(n)(n/s)=1$. This exact special case is also written
in \citet[\S3, pp.~4--5]{CrDarmon2009}. Grouping conjugate
generators $\alpha=(u+v\sqrt{-s})/2$ gives the integral formula
used by the coefficient stream:
\[
 b_n=\sum_{a^2=n,\ a>0}\left(\frac as\right)a
  +\sum_{\substack{u,v>0,\ u\equiv v\ (2)\\u^2+sv^2=4n}}
       \left(\frac{(s+1)u/2}{s}\right)u.
\]
For all four bases, exact PARI elliptic-curve coefficients equal
$b_n$ for $1\leq n\leq10000$.

This finite comparison proves equality of the entire modular forms.
For completeness, if a nonzero form $g$ of weight $k$ on
$\Gamma_0(N)$ has trivial character, take the product of
$g|_k\gamma$ over right coset representatives in
$\Gamma_0(N)\backslash\mathrm{SL}_2(\Z)$. It is a nonzero
holomorphic full-level form of weight $kI$, where
$I=N\prod_{p\mid N}(1+1/p)$. Every factor has nonnegative
order at infinity, so the product's order is at least
$\operatorname{ord}_\infty g$. The full-level valence formula
therefore implies $\operatorname{ord}_\infty g\leq kI/12$.
Equivalently, agreement through $\lfloor kI/12\rfloor$ proves
equality. The stronger arithmetic congruence statement is
\citet[Theorem 9.18]{CgzStein2007}; only the elementary
characteristic-zero argument is needed here. For $k=2,N=s^2$
the four bounds are
\[
 \left\lfloor\frac{s^2+s}{6}\right\rfloor
     =22,315,759,4455,
\]
all below $10000$. Applying the exact quadratic twist character
then identifies every coefficient in both the target and auxiliary
streams, including coefficients past the tested range. The exact
metadata transcript prints the four bounds and their comparisons.

\paragraph{5. Directed enclosures of both $L$-sums.}
For $c=2\pi/\sqrt N$, the root-number $+1$ formula is
\[
 L(E,1)=2\sum_{n\geq1}a_n\frac{e^{-cn}}n,
\]
while for root number $-1$ the first derivative is
\eqref{ari:eq:Cr-first}. The exact coefficients in the smaller cases
come from retained PARI files. The largest auxiliary stream uses
$s=163$, signed twist $-557403$, conductor
$8254937936042721$, and $M=867617471$; the target stream uses
$s=163,d=185801,M=289205824$. Both use the cutoff $cM\geq60$
up to rounding of the integer cutoff, with the actual $M$ used in
the tail bound. No precision conclusion depends on the nominal
cutoff.

Here are the finite-sum and rounding certificates implemented in
\path{code/B_C6_central_stream.cpp}. All intervals use 256-bit
MPFR with outward rounding at every operation. In a block centred
at $b$ with radius $R=128$, set
$S_k=\sum a_n(n-b)^k\in\Z$. For $g(z)=e^{-cz}/z$, put
\[
 C_0=1,\quad C_k=C_{k-1}c/k,\qquad
 T_0=1/b,\quad T_k=(T_{k-1}+C_k)/b.
\]
The coefficient of $(z-b)^k$ is $(-1)^ke^{-cb}T_k$.
For $h(z)=E_1(cz)/z$, put
\[
 e_0=E_1(cb),\quad e_k=(-1)^ke^{-cb}T_{k-1}/k\ (k\geq1),
 \qquad h_0=e_0/b,\quad h_k=(e_k-h_{k-1})/b.
\]
The block sums are the finite sums of these interval coefficients
times the exact integers $S_k$.

On $|z-b|=b/2$, the two kernels have bounds
\[
 |g(z)|\leq\frac{2e^{-cb/2}}b=:U_g,
 \qquad
 |h(z)|\leq\frac{4e^{-cb/2}}{cb^2}=:U_h.
\]
For the second bound, use
$E_1(cz)=\int_1^\infty e^{-czt}\,dt/t$ when $\Re z>0$.
Thus truncation at degree $k$ contributes at most
\[
 U_g\frac{(2R/b)^{k+1}}{1-2R/b}\sum|a_n|
 \quad\hbox{or}\quad
 U_h\frac{(2R/b)^{k+1}}{1-2R/b}\sum|a_n|.
\]
The degree is selected to make the geometric factor small, but
the actual remainder is separately evaluated with upward rounding.
For the derivative stream, the initial $E_1$ in each segment is
evaluated using correctly rounded $-\operatorname{Ei}(-x)$ and
monotonicity. Subsequent centres, spaced by $W=257$, use
\[
 E_1(c(b+W))=E_1(cb)-\int_0^Wg(b+x)\,dx.
\]
Integrating the same Taylor expansion gives the interval for this
advance. Its omitted part is bounded by
$WU_g(2W/b)^{k+1}/(1-2W/b)$ and added to both endpoints.
The first $10280$ terms are summed directly. Segment boundaries
are multiples of $257$, so no coefficient is omitted or repeated.

Hasse's bound and multiplicativity give $|a_n|\leq2n$.
The infinite tails, including the factor two, are consequently at
most
\[
 \frac{4e^{-c(M+1)}}{1-e^{-c}},\qquad
 \frac{4e^{-c(M+1)}}{c(M+1)(1-e^{-c})},
\]
respectively. The program evaluates both with the downward bound
for $c$ and upward bounds for the remaining positive expressions.
It finally adds the block remainders, interval arithmetic errors,
and infinite tail. For row eleven the central-value tail is below
$5.07\cdot10^{-19}$, and the derivative tail is below
$2.82\cdot10^{-21}$. These modest precisions already suffice for
rational separation. All twelve completed interval outputs are
retained, as are the exact coefficient files.

\paragraph{6. Periods, heights, and exact rational separation.}
The periods are enclosed independently with 512-bit Arb. Let
$\tau=i\sqrt s$ for $i=2,4$ and
$\tau=(1+i\sqrt s)/2$ for $i=6,9,10,11$; here
$s=3,7,11,43,67,163$, respectively. The exact CM identification
and reduced-form class-number check are part of C.6.5. Put
$T=\Im\tau$, $q_\tau=e^{2\pi i\tau}$, and
$r=|q_\tau|<10^{-4}$. We enclose
\[
 E_4(\tau)=1+240\sum_{n=1}^{100}\sigma_3(n)q_\tau^n
                   +\text{tail}.
\]
Since $\sigma_3(n)<2n^3$, with $A=101$ the absolute tail is at most
\[
 480r^A\left\{
 \frac{A^3}{1-r}+\frac{3A^2r}{(1-r)^2}
 +\frac{3Ar(1+r)}{(1-r)^3}
 +\frac{r(1+4r+r^2)}{(1-r)^4}\right\}.
\]
The exact invariants $c_4,c_6$ of each minimal equation give
$\kappa=2\pi(E_4(\tau)/c_4)^{1/4}$. The sign of $E_6(\tau)$
is positive: $\sigma_5(n)<2n^5$ and
$n^5\leq32^{n-1}$ give
$|E_6(\tau)-1|\leq1008r/(1-32r)<1$.
Thus the full real period is
\[
 \Omega=n_{\R}\kappa\begin{cases}
 1,&c_6>0,\\
 T,&c_6<0,\ \Delta>0,\\
 2T,&c_6<0,\ \Delta<0,
 \end{cases}
 \qquad n_{\R}=\begin{cases}2,&\Delta>0,\\1,&\Delta<0.\end{cases}
\]
These formulas apply separately to the target and auxiliary curve.
Every operation, the Fourier tail, and the fourth root is enclosed
by Arb. The height intervals are the independent certificates
\path{records/C_cm_height_arb.json} of C.6.5.

The final computation is
\path{code/B_C6_exact_factors.py}. First, the intervals for
$r_i$ isolate a unique integer after multiplication by the exact
$t_bn_b$ in \eqref{ari:eq:Cgz-central-integrality}. This proves the
six auxiliary rational values, with no BSD input. Next form the
interval $L'(E_i,1)/(\Omega_i\widehat h_O(P_i))$. The enclosures
below are deliberately enlarged decimal summaries of the recorded
Arb intervals:
\[
\begin{array}{c|c|c}
 i&|r_i-r_i^{\rm listed}|\text{ is at most}
   &|q_i-q_i^{\rm listed}|\text{ is at most}\\\hline
 2&4\cdot10^{-22}&2\cdot10^{-53}\\
 4&4\cdot10^{-20}&3\cdot10^{-45}\\
 6&4\cdot10^{-22}&2\cdot10^{-53}\\
 9&5\cdot10^{-21}&4\cdot10^{-53}\\
 10&3\cdot10^{-20}&10^{-23}\\
 11&7\cdot10^{-15}&8\cdot10^{-19}
\end{array}
\]
Each last-column bound is strictly less than $1/B_i$. A rational
number with denominator at most $B_i$ differing from the listed
integer would differ by at least $1/B_i$. Equation
\eqref{ari:eq:Cgz-denominator} therefore makes every listed $q_i$
exact. The same equation gives the listed $|a_i|$ by integer
square-root checks. The JSON certificate records all input
intervals, exact invariants, rational values, and positive separation
margins. Its stdout concludes that all six remaining factors are
exactly certified. This completes the proof.
\end{proof}
